% BHU PYQ -- Professional Exam Template
\documentclass[11pt,a4paper]{article}

\usepackage[a4paper,top=2.5cm,bottom=2.5cm,left=2.8cm,right=2.8cm,headheight=14pt]{geometry}
\usepackage{amsmath,amssymb}
\usepackage[T1]{fontenc}
\usepackage[utf8]{inputenc}
\usepackage{lmodern}
\usepackage{microtype}
\usepackage{fancyhdr}
\usepackage{enumitem}
\usepackage{listings}
\usepackage{hyperref}
\usepackage{lastpage}
\usepackage{parskip}

\hypersetup{colorlinks=true,urlcolor=black,linkcolor=black,
  pdftitle={CS-103: Numerical Computing -- BHU PYQ}}

\pagestyle{fancy}
\fancyhf{}
\renewcommand{\headrulewidth}{0.4pt}
\renewcommand{\footrulewidth}{0.4pt}
\fancyhead[L]{\small   \textit{Banaras Hindu University}}
\fancyhead[R]{\small   \textit{CS-103: Numerical Computing}}
\fancyfoot[C]{\small Page  \thepage\ of \pageref{LastPage}}

\lstset{
  basicstyle=\small tfamily,
  frame=single,
  breaklines=true,
  columns=flexible,
  keepspaces=true
}


\newlist{parts}{enumerate}{2}
\setlist[parts,1]{label=(\alph*),leftmargin=*,itemsep=4pt,topsep=3pt}
\setlist[parts,2]{label=(\roman*),leftmargin=*,itemsep=2pt,topsep=2pt}

\newcommand{\pts}[1]{\hfill   \textit{[#1]}}


\begin{document}

\begin{center}
  {\large\bfseries BANARAS HINDU UNIVERSITY}\\[2pt]
  {
\normalsize B.Sc. (Hons.) Semester III Examination, 2022-23}\\[6pt]
  \rule{0.8\linewidth}{0.5pt}\\[6pt]
  {\large\bfseries Paper: CS-103 --- Numerical Computing}\\[4pt]
  {
\normalsize Time: Three hours \hfill Maximum Marks: 70}
\end{center}

\rule{\linewidth}{0.4pt}

\smallskip



\noindent   \textit{Instructions: ATTEMPT ANY FIVE QUESTIONS. THE FIGURES IN THE RIGHT-HAND MARGIN INDICATE}

\bigskip

\medskip


\noindent   \textbf{Question 1.}

\begin{parts}
  \item Express the following decimal numbers in binary, octal and hexadecimal:\pts{6}
  
\begin{parts}
    \item 22.623
    \item -10.215
  \end{parts}
  \item What are the three different types of errors encountered in numerical calculations? Explain each of them briefly.\pts{3}
  \item Determine the roots of the following equation using Newton-Raphson method correct to three decimal places:\pts{5}
  \[ x^3 - 5x + 3 = 0 \]
\end{parts}

\medskip


\noindent   \textbf{Question 2.}

\begin{parts}
  \item Derive the Regula-Falsi method for calculation of roots of non-linear equation. Write the steps of calculation of roots using this method. Discuss the scenario where this method would give large errors.\pts{8}
  \item Calculate the roots of the following equation using Regula-Falsi method, correct to three decimal places:\pts{6}
  \[ x^3 + x^2 + x + 7 = 0 \]
\end{parts}

\medskip


\noindent   \textbf{Question 3.}

\begin{parts}
  \item Describe the LU decomposition method for solving system of linear equations.\pts{2}
  \item Use the LU decomposition method to solve the following system of linear equations:\pts{7}
  \[
  
\begin{aligned}
    5x - 2y + z &= 4 \\
    7x + y - 5z &= 8 \\
    3x + 7y + 4z &= 10
  \end{aligned}
  \]
  \item Describe the method of diagonalization of a real symmetric matrix using Jacobi method.\pts{5}
\end{parts}

\medskip


\noindent   \textbf{Question 4.} For the matrix $A$ given below, calculate:
\[ A = 
\begin{bmatrix} 4 & 3 & 2 \\ 2 & 3 & 4 \\ 1 & 2 & 1 \end{bmatrix} \]

\begin{parts}
  \item The largest Eigen value and the corresponding Eigen vector.\pts{7}
  \item Inverse of the matrix.\pts{7}
\end{parts}

\medskip


\noindent   \textbf{Question 5.}

\begin{parts}
  \item Derive Simpson's $1/3^{ \text{rd}}$ Rule for numerical integration.\pts{4}
  \item Integrate the following function using:
  \[ I = \int_{0}^{1} \frac{1}{1+x} \, dx \]
  
\begin{parts}
    \item Trapezoidal Rule
    \item Simpson's $1/3^{ \text{rd}}$ Rule
  \end{parts}
  using $h = 0.125$. Also compare the results obtained in each case with the actual value of the integral.\pts{10}
\end{parts}

\medskip


\noindent   \textbf{Question 6.}

\begin{parts}
  \item Derive the Newton-Gregory forward interpolation formula.\pts{4}
  \item Use Newton's Divided difference method to calculate $f(0.25)$ and $f(0.27)$ for the following tabulated function:\pts{8}
  
\begin{center}
  
\begin{tabular}{|c|c|c|c|c|c|c|}
    \hline
    $x$ & 0.20 & 0.22 & 0.24 & 0.26 & 0.28 & 0.30 \\
    \hline
    $f(x)$ & 1.6596 & 1.6698 & 1.6804 & 1.6912 & 1.7024 & 1.7139 \\
    \hline
  \end{tabular}
  \end{center}
  \item State two applications of interpolation.\pts{2}
\end{parts}

\medskip


\noindent   \textbf{Question 7.}

\begin{parts}
  \item Use Picard's method to solve the differential equation:
  \[ \frac{dy}{dx} = 1 + xy, \quad y(0) = 1 \]
  Find the value of $y$ for $x = 0.5$ correct to 3 decimal places.\pts{7}
  \item Use Predictor-Corrector method to solve the following differential equation:
  \[ \frac{dy}{dx} = y(1 + x^2), \quad y(0) = 1 \]
  Use $h = 0.05$ and compute $y(0.3)$ correct to 3 decimal places.\pts{7}
\end{parts}

\vfill
\rule{\linewidth}{0.4pt}

\begin{center}\small
\href{https://bhu-exam.vercel.app/}{Click here to visit our website}\\[4pt]\href{https://me-aryan.vercel.app/}{Click here to visit our contributor}\\[4pt]\href{https://quashan-me.vercel.app/}{Click here to visit our contributor}
\end{center}
\end{document}